\documentclass[11pt]{article}
\newcommand{\DocShort}{Anexo SWMM}
\usepackage{fontspec}
\usepackage[spanish,es-noshorthands]{babel}
\decimalpoint
\usepackage{geometry}
\geometry{letterpaper, top=2.5cm, bottom=2.5cm, left=2.8cm, right=2.8cm}
\usepackage[expansion=false]{microtype}
\usepackage{parskip}
\setlength{\parskip}{6pt}
\usepackage{amsmath,amssymb}
\usepackage{booktabs}
\usepackage{array}
\usepackage{longtable}
\usepackage{caption}
\usepackage[dvipsnames,table]{xcolor}
\usepackage{enumitem}
\setlist[itemize]{topsep=2pt, itemsep=2pt, parsep=0pt}
\setlist[enumerate]{topsep=2pt, itemsep=2pt, parsep=0pt}
\usepackage{fancyhdr}
\usepackage[hidelinks,pdfencoding=auto]{hyperref}
\usepackage{titlesec}
\usepackage{tcolorbox}
\tcbuselibrary{breakable}
\usepackage{pdfpages}

\definecolor{darkblue}{HTML}{1B3A5C}
\definecolor{medblue}{HTML}{2E6DA4}
\definecolor{lightblue}{HTML}{D6E8F7}
\definecolor{teal}{HTML}{1A7A6E}
\definecolor{lightteal}{HTML}{D0EDEA}
\definecolor{lightgrey}{HTML}{F5F5F5}
\definecolor{midgrey}{HTML}{6F6F6F}
\definecolor{warnyellow}{HTML}{FFF3CD}

\titleformat{\section}
{\large\bfseries\color{darkblue}}
{\thesection}{0.75em}{}[{\color{medblue}\titlerule[0.8pt]}]
\titleformat{\subsection}
{\normalsize\bfseries\color{darkblue}}
{\thesubsection}{0.75em}{}
\titleformat{\subsubsection}
{\normalsize\bfseries\color{medblue}}
{\thesubsubsection}{0.75em}{}

\pagestyle{fancy}
\fancyhf{}
\fancyhead[L]{\small\color{midgrey}ICV 442 Acueductos y Alcantarillados}
\fancyhead[R]{\small\color{midgrey}\DocShort}
\fancyfoot[C]{\small\color{midgrey}\thepage}
\renewcommand{\headrulewidth}{0.4pt}
\renewcommand{\footrulewidth}{0pt}

\rowcolors{2}{lightgrey}{white}
\renewcommand{\arraystretch}{1.18}

\newtcolorbox{infobox}{
  breakable,
  colback=lightblue,
  colframe=medblue,
  boxrule=0.8pt,
  arc=3pt,
  left=8pt,
  right=8pt,
  top=7pt,
  bottom=7pt
}

\newtcolorbox{warnbox}{
  breakable,
  colback=warnyellow,
  colframe=orange!70!black,
  boxrule=0.8pt,
  arc=3pt,
  left=8pt,
  right=8pt,
  top=7pt,
  bottom=7pt
}

\newcommand{\doccover}[3]{
  \begin{center}
  {\small\color{midgrey}Pontificia Universidad Cat\'olica Madre y Maestra\\Escuela de Ingenier\'ia Civil y Ambiental}

  \vspace{1.0cm}
  {\color{medblue}\rule{\linewidth}{2pt}}
  \vspace{0.3cm}

  {\Large\bfseries\color{darkblue}ICV-442-T Acueductos y Alcantarillado}\\[0.35em]
  {\LARGE\bfseries\color{darkblue}#1}\\[0.25em]
  {\large\itshape #2}

  \vspace{0.3cm}
  {\color{medblue}\rule{\linewidth}{2pt}}
  \vspace{0.9cm}
  \end{center}

  \begin{center}
  \begin{tabular}{ll}
  \toprule
  \textbf{Asignatura} & ICV-442-T Acueductos y Alcantarillado \\
  \textbf{Profesor} & Ricardo Rafael Hern\'andez Moreira, PhD \\
  \textbf{Periodo acad\'emico} & Mayo--agosto de 2026 \\
  \textbf{Proyecto} & Dise\~no de sistemas de saneamiento para Las Charcas, Azua \\
  \textbf{Versi\'on} & #3 \\
  \bottomrule
  \end{tabular}
  \end{center}

  \vspace{0.4cm}
}


\begin{document}
\doccover{Anexo: patr\'on de descarga para SWMM}{Alcantarillado sanitario}{18 de mayo de 2026}

\section{Uso del patr\'on}
El patr\'on siguiente se propone para representar la variaci\'on horaria del caudal sanitario de tiempo seco. Su funci\'on es operativa: apoyar la configuraci\'on del modelo. El equipo sigue siendo responsable de justificar los aportes sanitarios y los factores empleados.

\section{Tabla horaria}
\begin{center}
\begin{tabular}{cccccccc}
\toprule
Hora & Mult. & Hora & Mult. & Hora & Mult. & Hora & Mult. \\
\midrule
00:00 & 0.35 & 06:00 & 0.50 & 12:00 & 0.85 & 18:00 & 0.95 \\
01:00 & 0.30 & 07:00 & 0.75 & 13:00 & 0.80 & 19:00 & 0.95 \\
02:00 & 0.25 & 08:00 & 0.95 & 14:00 & 0.75 & 20:00 & 1.00 \\
03:00 & 0.25 & 09:00 & 1.00 & 15:00 & 0.70 & 21:00 & 0.85 \\
04:00 & 0.25 & 10:00 & 0.90 & 16:00 & 0.75 & 22:00 & 0.70 \\
05:00 & 0.35 & 11:00 & 0.80 & 17:00 & 0.85 & 23:00 & 0.50 \\
\bottomrule
\end{tabular}
\end{center}

\section{Bloque para archivo \texttt{.inp}}
\begin{verbatim}
[PATTERNS]
DescargaSanitaria_LasCharcas    HOURLY
    1 0.35   2 0.30   3 0.25   4 0.25   5 0.25   6 0.35
    7 0.50   8 0.75   9 0.95  10 1.00  11 0.90  12 0.80
   13 0.85  14 0.80  15 0.75  16 0.70  17 0.75  18 0.85
   19 0.95  20 0.95  21 1.00  22 0.85  23 0.70  24 0.50
\end{verbatim}

\section{Nota de uso}
\begin{itemize}
\item Este patr\'on se aplica a caudales sanitarios de tiempo seco.
\item No es necesario forzar la documentaci\'on con discusiones sobre m\'etodos de infiltraci\'on de lluvia ajenos al problema sanitario separativo.
\item Si el equipo incorpora aportes adicionales constantes, debe explicarlos por separado en la memoria.
\end{itemize}

\end{document}
